% =========================================================================
% Figura (apuntes de teoría, pág. II-3): decaimiento 1/r^2 frente a 1/r^3.
% Se cortan en r = 1: más cerca domina el término más rápido; más lejos,
% el más lento.
% =========================================================================
\begin{figure}[ht]
\centering
\begin{tikzpicture}
\begin{axis}[tema, width=8.6cm, height=5.2cm, xmin=0.4, xmax=4, ymin=0, ymax=4.2,
             xlabel={$r$ (u.\,arb.)}, ylabel={campo (u.\,arb.)}, grid=major,
             xtick={0.5,1,...,4}, ytick={0,1,...,4}, clip=true,
             legend pos=north east, legend cell align=left]
  \addplot[mmaazul, line width=1.1pt, domain=0.45:4, samples=150] {1/x^2};
  \addlegendentry{$1/r^{2}$ (carga)}
  \addplot[mmanaranja!90!black, line width=1.1pt, domain=0.6:4, samples=150] {1/x^3};
  \addlegendentry{$1/r^{3}$ (dipolo)}
  \draw[punteado] (axis cs:1,0) -- (axis cs:1,1);
  \node[font=\scriptsize, text=grispizarra, anchor=south west] at (axis cs:1.05,1.05) {$r=1$};
\end{axis}
\end{tikzpicture}
\caption{Comparación del decaimiento de una contribución monopolar ($1/r^{2}$, campo de una carga) y otra dipolar ($1/r^{3}$, campo de
un dipolo) con la misma intensidad en $r=1$. Para $r<1$ domina el término dipolar; para $r>1$, el de la carga, y la diferencia se agranda
rápidamente con la distancia.}
\label{fig:potencias}
\end{figure}
